\section{Ce qui reste à démontrer}\label{sec:reste}

Le tableau~\ref{tab:reste} donne, pour chaque cote, les énoncés non démontrés ici : ceux que visait le choix de la cote, puis les autres énoncés de la lecture.

\begingroup\small
\begin{longtable}{>{\raggedright\arraybackslash}p{1.2cm}>{\raggedright\arraybackslash}p{6.4cm}>{\raggedright\arraybackslash}p{6.4cm}}
\caption{Énoncés des treize cotes non démontrés ici.}\label{tab:reste}\\\toprule
Cote & Énoncés visés, non démontrés & Autres énoncés de la lecture\\\midrule\endfirsthead
\toprule Cote & Énoncés visés, non démontrés & Autres énoncés de la lecture\\\midrule\endhead
42 & Lemme~1 (le lieu de cohérence est ouvert) ; corollaire~1 ($\Gamma(X,\mathcal O_X)$ comme limite des germes). & La proposition de la p.~3 ; le lemme~3 et les corollaires~2 et~3 (pleine fidélité) ; les deux contre-exemples ; la forme $(A_f)_{\mathfrak q} \to A_{\mathfrak p}$ de~(2). \\
153 & --- & Les co-opérations et leur unicité ; Pólya et Vandermonde ; $\Omega = \Omega_0 \oplus \Omega^0$ ; les exemples et contre-exemples des pp.~7 et~9. \\
22 & Lissité, non-ramification et étalitude formelles, et leurs critères de relèvement. & Dans~3.5, la réduction à un système indexé par $\mathbb N$ ; 3.2 et 3.6--3.13 (algèbres de Cohen, homomorphismes réguliers) ; pp.~1--8 et~18--28. \\
21 & Profondeur le long d'un fermé ; prolongement de type Hartogs. & 2.1--2.14 (Eilenberg--Watts, module représentant, premiers associés) ; 2.15 comme équivalence de catégories ; 2.17 ; pp.~15--28 (faisceaux, catégories fibrées, torseurs). \\
39 & La préparation (unité $\times$ polynôme distingué) ; paires henséliennes ; séries algébriques. & Le cas d'une topologie linéaire quelconque ; le cas $I$-adique sous $a_i \in \sqrt I$ ; l'axiomatique des pp.~2--16. \\
125 & Profondeur le long d'un fermé ; prolongement de type Hartogs. & La dualité de Serre (admise) ; les étages~4 et~5 (complétion formelle) ; les annulations pour $\mathcal O(n)$ ; la dimension cohomologique de $P - X$ ; les Ext des pp.~4--5 ; III-5 (Koszul) ; III-18 pour les faisceaux. \\
41 & Constructibilité ; ouverture universelle ; irréductibilité géométrique des fibres. & Le cas \q{$Y$ localement noethérien} ; la fibre schématique comme $f^{-1}(y)$ ; les définitions~2 et~3 ; la représentabilité des sections monogènes ; le théorème des pp.~8--9. \\
113 & La forme $k$-linéaire de l'invariance de Morita. & L'indépendance de la structure $k$-additive ; $\mathrm{Kar}\,P \simeq \mathrm{Proj}_{\mathrm{tf}}(P^k)$ ; Freyd--Mitchell ; l'énoncé 2-fonctoriel ; pp.~5--12 (produit tensoriel, contre-exemple~($\circledast$), produit dérivé). \\
161-1 & La préservation des colimites par l'adjonction idempotente. & La condition duale~(2) ; les cinq conditions de la p.~3 ; $E_0 \simeq F_0$ ; les conditions 1), 3), 4) de la p.~5 ; la descente de la p.~6 ; pp.~9--19. \\
104 & Les identités simpliciales ; la bonne fondation (catégories de Reedy). & $M \simeq M_0$ et la réciproque (pp.~9--10) ; sphéricité et régularité de $|I_n^*|$ ; la p.~11 ; pp.~12--16 ; les conditions $(C_i)$ de la p.~4 ; les sous-objets des exemples. \\
158 & --- & Le groupe fondamental (pp.~19--20) ; la catégorie $0$-connexe non $1$-connexe (pp.~21--41) ; la proposition~2 \q{(??)} et le corollaire \q{Faux !} ; pp.~50--53 ; $\mathrm{Ep}(E)$ et $\mathrm{Mon}(E)$ ; pp.~60--83. \\
152 & Plongements dans une surface orientée et circularisations (p.~4). & Le critère encadré des arbres (p.~9) ; la fonctorialité ; pp.~2, 3, 10--12. \\
88 & --- (la présentation de Coxeter est une définition). & $\mathrm{Aut}(\Sigma) \to \mathrm{Aut}(\Phi(\Sigma))$ (pp.~3--4) ; les autres relations entre les $\sigma_i$ ; la triangulation de la sphère ; les articles (1)--(5) et~(7) ; le contenu géométrique de~(6). \\\bottomrule
\end{longtable}
\endgroup

La seule erreur trouvée dans un énoncé des lectures (cote~158) portait sur une glose moderne, non sur ce que la page énonce. Ce qui reste attend pour l'essentiel des théories classiques : surfaces, dualité de Serre, complétion formelle, groupe fondamental d'une catégorie, lissité formelle.
